\chapter{Jalon M4 --- Restitution BI (Grafana)}
\label{chap:m4}

M4 n'est pas un bloc de compétences RNCP mais un jalon transverse,
ajouté en cours de programme (24/09/2026) après un retour de formateur :
les livrables produits pour C16 (vues SQL + notebook) et C20 (catalogue
par introspection) suffisaient aux critères déjà évalués, mais pas
comme démonstration de restitution/monitoring dans un outil de BI
dédié. Clos le 19/10/2026, les 2 issues du jalon sont vérifiées en
conditions réelles -- détail exhaustif dans
\texttt{docs/comptes\_rendus/M4.md}.

\section{Choix de l'outil}

Grafana est retenu, auto-hébergé, plutôt que Power~BI (nécessite le
Power~BI Service, cloud) ou Databricks~SQL/Lakeview (nécessite un
\textit{workspace} Databricks) -- cohérent avec l'infrastructure 100~\%
locale du reste du programme, décidé avant tout codage. Le notebook SLA
existant n'est ni retiré ni réécrit : il reste la preuve exploratoire
telle que livrée à la clôture du Bloc~3, Grafana devenant l'outil de
restitution \textit{de production} en complément.

\section{C16bis --- Restitution de l'entrepôt OMEGA BI}

Les vues SQL déjà exposées pour C16 sont restituées dans un dashboard
Grafana dédié, connecté en lecture seule via le rôle \texttt{bi\_reader}
déjà scopé en C14. Trois bugs réels ont été trouvés uniquement en
testant depuis un vrai navigateur (champ requis manquant au provisioning
de la source de données, mode d'édition de requête non forcé sur les
panels, agrégation silencieuse masquant deux clients sur trois sur un
panel) -- tous invisibles depuis les vérifications par appel direct à
l'API de Grafana, qui passaient toutes. Généralisé en leçon
méthodologique : une vérification par API seule ne prouve jamais
qu'un rendu côté client fonctionne pour l'utilisateur final.

\section{C20bis --- Monitorage du \textit{data lake}}

Une première conception du monitorage du lake, limitée au suivi de la
fraîcheur des données, s'est révélée incomplète après lecture du
référentiel de compétences complet (copié dans le dépôt,
\texttt{docs/reference/}) : celui-ci exige explicitement le suivi des
« conditions matérielles et applicatives », et qu'une alerte soit
« générée lors d'une rupture de service ». Corrigé par l'ajout d'un
relevé de santé de l'infrastructure (MinIO accessible, capacité
utilisée) et d'une règle d'alerte Grafana réelle.

La chaîne d'alerte relie une règle Grafana, un service interne
authentifié, et un relais d'e-mail réel via la bibliothèque
\texttt{smtplib} vers un capteur SMTP local (MailHog) -- jamais un
vrai fournisseur externe, pour rester démonstrable sans identifiant
réel à gérer ni envoi parasite à chaque exécution automatisée des
tests. Testée de bout en bout, deux fois : une coupure réelle de MinIO
déclenche un relevé en échec, l'alerte Grafana passe à l'état actif en
une trentaine de secondes, et un e-mail réel est reçu dans MailHog ;
la remise en service résout l'alerte. Une relecture de code a ensuite
mis au jour deux trous de robustesse (un corps de requête mal formé
faisait échouer le service sans réponse exploitable, un échec
d'envoi d'e-mail n'était pas distingué d'une autre erreur) et l'absence
d'authentification sur ce service, exposé sur l'hôte -- les trois
corrigés et revérifiés avant la clôture du jalon.

\section{Portée budgétaire}

Comme les jalons transverses de rédaction du rapport et de préparation
de la soutenance, M4 n'a jamais eu de ligne dédiée dans le budget
prévisionnel initial (180~000~€, voir chapitre~\ref{chap:bloc1}) --
staffing interne uniquement, les outils retenus (Grafana, MailHog)
étant open source et auto-hébergés. Aucun dépassement constaté.
